% !TeX program = xelatex
% UTF-8 transcription of Google Drive Notebook/neuroblastoma, source 158.
% pic01.jpg: 20261013_065004465.jpg (handwritten page 1).
% pic02.jpg: 20261013_065009739.jpg (handwritten page 2).
% pic03.jpg: 20261013_065016278.jpg (handwritten page 3).
% pic04.jpg: 20261013_065020482.jpg (handwritten page 4).
% Source statements and numbers are transcribed without medical correction.
% Compile directly with XeLaTeX, or input into a XePersian document.
\newif\ifHemNeuroStandalone
\ifdefined\settextfont
  \HemNeuroStandalonefalse
\else
  \HemNeuroStandalonetrue
\fi
\ifHemNeuroStandalone
  \documentclass[12pt,a4paper]{article}
  \usepackage[margin=20mm]{geometry}
  \usepackage{amsmath,amssymb,xcolor}
  \usepackage{xepersian}
  \IfFontExistsTF{Amiri}{\settextfont{Amiri}}{\settextfont{FreeSerif}}
  \setlatintextfont{FreeSerif}
  \begin{document}
\fi
\begingroup
\setlength{\parindent}{0pt}
\setlength{\parskip}{4pt}
\emergencystretch=3em
\newcommand{\HemNeuroEn}[1]{\lr{#1}}
\ifdefined\colorbox
  \newcommand{\HemNeuroHl}[1]{\colorbox{yellow!35}{\strut #1}}
\else
  \newcommand{\HemNeuroHl}[1]{\textbf{#1}}
\fi

\section*{۱۵۸ ـ نوروبلاستوما؛ صفحهٔ ۱}
\subsection*{تومورهای نوروبلاستیک و منشأ}
\begin{itemize}
\item تومورهای نوروبلاستیک:
  \begin{itemize}
  \item نوروبلاستوما.
  \item گانگلیونوروبلاستوما.
  \item گانگلیونوروما.
  \end{itemize}
\item منشأ: نورال کرست؛ تمایز به \HemNeuroHl{مدولا آدرنال} و \HemNeuroHl{زنجیرهٔ سمپاتیک}.
\item فئوکروموسایتوما و پاراگانگلیوما از سلول‌های کرومافین منشأ می‌گیرند:
  \HemNeuroEn{Migrates from Neural Crest}.
\end{itemize}

\subsection*{اپیدمیولوژی و عوامل مادری}
\begin{itemize}
\item سومین سرطان شایع در کودکان است؛ پس از:
  \begin{enumerate}
  \item لوکمی.
  \item تومور مغز.
  \end{enumerate}
\item در گروه شیرخواران، در سفیدپوست‌ها شایع‌تر است، ولی در بقیهٔ کودکان در دو گروه برابر است.
\item در دو جنس برابر است؛ در پسرها اندکی بیشتر.
\item ریسک‌فاکتورهای مادری نوروبلاستوما:
  \begin{itemize}
  \item مصرف اوپیوئید.
  \item کمبود فولات.
  \end{itemize}
\item اکثر موارد اسپورادیک هستند.
\end{itemize}

\clearpage
\section*{۱۵۸ ـ نوروبلاستوما؛ صفحهٔ ۲}
\subsection*{محل اولیه و متاستاز}
\begin{itemize}
\item شایع‌ترین محل: \HemNeuroHl{غدهٔ آدرنال} (۴۰\%)، سپس \HemNeuroHl{شکم} (۲۵\%).
\item ۵۰\% بیماران در زمان تشخیص متاستاز دارند.
\item محل‌های شایع متاستاز:
  \begin{itemize}
  \item \HemNeuroHl{لنف‌نود}.
  \item مغز استخوان.
  \item جمجمه و دورا: هایپرکلسمی.
  \item چشم، کبد، پوست: \HemNeuroHl{بلوبری مافین ندول} و اکیموز پری‌اوربیتال.
  \end{itemize}
\end{itemize}

\subsection*{علائم و نشانه‌ها}
\begin{itemize}
\item تودهٔ شکمی.
\item درد شکم یا پوست.
\item اکیموز پری‌اوربیتال (\HemNeuroHl{راکون آی}): متاستاز به چشم.
\item \HemNeuroHl{سندرم هورنر}: \HemNeuroHl{میوز، پتوز، انوفتالموس}.
\item درد کمر، ضعف عضلات: \HemNeuroHl{\HemNeuroEn{Cord compression}}.
\item \HemNeuroEn{HTN}.
\item \HemNeuroHl{سندرم اپسوکلونوس مایوکلونوس}: پارانئوپلاستیک.
\item \HemNeuroEn{Palpable nontender subcutaneous nodules} (\HemNeuroHl{بلوبری مافین}).
\end{itemize}

\subsection*{معاینه}
\begin{itemize}
\item تودهٔ شکمی معمولاً \HemNeuroHl{از خط وسط عبور می‌کند} (برخلاف تومور ویلمز).
\item در معاینهٔ شکم کودک اگر توده دیدیم، \HemNeuroHl{دستکاری نمی‌کنیم}.
\item در معاینه، توده معمولاً:
  \begin{itemize}
  \item \HemNeuroEn{Non-tender}.
  \item \HemNeuroEn{Fixed}.
  \item \HemNeuroEn{Firm}.
  \end{itemize}
\end{itemize}

\clearpage
\section*{۱۵۸ ـ نوروبلاستوما؛ صفحهٔ ۳}
\subsection*{سندرم‌های پارانئوپلاستیک}
\begin{itemize}
\item \HemNeuroEn{Opsoclonus myoclonus} (\HemNeuroHl{\HemNeuroEn{OMS}}):
  \begin{itemize}
  \item در ۱ تا ۳\% نوروبلاستوماها.
  \item پاتوژنز \HemNeuroHl{خودایمن} دارد.
  \item ۵۰\% موارد \HemNeuroEn{OMS} نوروبلاستوما دارند.
  \item \HemNeuroHl{\HemNeuroEn{Rapid dancing eye movements}}.
  \item \HemNeuroHl{\HemNeuroEn{Rhythmic jerking}}؛
    \HemNeuroEn{(myoclonus) involving limbs or trunk}.
  \item \HemNeuroEn{and/or} \HemNeuroHl{\HemNeuroEn{ataxia}}.
  \end{itemize}
\item \HemNeuroEn{Secretion of VIP}:
  \begin{itemize}
  \item دیستنشن و درد شکمی.
  \item اسهال ترشحی و هایپوکالمی.
  \item \HemNeuroHl{\HemNeuroEn{Serum VIP \(>75\) pg/ml}}.
  \end{itemize}
\end{itemize}

\subsection*{آزمایشگاه در نوروبلاستوما}
در ۸۰\% نوروبلاستوماها افزایش یافته است:
\begin{itemize}
\item \HemNeuroEn{Serum vanillylmandelic acid} (\HemNeuroHl{\HemNeuroEn{VMA}}).
\item \HemNeuroEn{Serum hemovanillic acid} (\HemNeuroHl{\HemNeuroEn{HVA}}).
% The source writes "hemovanillic"; the spelling is retained.
\item \HemNeuroEn{Urine HVA, VMA}.
\end{itemize}

\subsection*{سن بیماران}
متوسط سن بیماران \HemNeuroHl{۱۸ ماهگی} است (البته تا ۱۷ سالگی).

\clearpage
\section*{۱۵۸ ـ نوروبلاستوما؛ صفحهٔ ۴}
\subsection*{تصویربرداری و بررسی مغز استخوان}
\begin{itemize}
\item تصویربرداری‌ها: \HemNeuroHl{\HemNeuroEn{MIBG Scan}}، \HemNeuroEn{MRI or CT}.
\item باید تحت بیوپسی / آسپیراسیون \HemNeuroHl{دوطرفهٔ} مغز استخوان قرار بگیرند.
\end{itemize}

\subsection*{درمان}
\begin{itemize}
\item قدم اولیه: جراحی و خارج کردن تمام تومور.
\item ریسک کم: صرفاً جراحی.
\item ریسک متوسط: جراحی + ۴ تا ۸ دورهٔ \HemNeuroEn{ChT}.
\end{itemize}

\endgroup
\ifHemNeuroStandalone
  \def\HemNeuroFinish{\end{document}}
\else
  \let\HemNeuroFinish\relax
\fi
\HemNeuroFinish
